% !TEX spellcheck = en_US



% ~~~~~~~~~~~~~~~~~~~~~~~~~~~~~~~~~~~~~~~ V
% (NC) 2026 Haluk Bingol
% github.com/halukbingol/LaTeX-Templates
%
% Licensees may copy, distribute, display, and perform the work and make derivative works 
% and remixes based on it only for non-commercial purposes. 
% ~~~~~~~~~~~~~~~~~~~~~~~~~~~~~~~~~~~~~~~ A




% =====================================================================
%  Linked Lists -- Sophomore Level
%  ---------------------------------------------------------------------
%  Built on hblib/Template/hblibTemplate.tex with the libraries in ../../hblib/.
%  Programs and their real outputs are in codeoutput/:
%     codeoutput/<Name>.java   a program (classes without main() are only compiled)
%     codeoutput/<Name>.txt    what it printed, regenerated by:  bash run_examples.sh
% =====================================================================

\documentclass[aspectratio=169,10pt,compress]{beamer}

% ~~~~~~~~~~~~~~~~~~~~~~~~~~~~~~~~~~~~~~~~~~~~~~~~~~~~~~~~~~~~~~~~~ V hblib
	\usepackage{../../hblib/hblibPresentation} % lib presentation: theme, i/N, \hSection, algorithm2e, ...
	\usepackage{../../hblib/hblibCodeoutput} % lib code + output: \lstcode, \lstterminal, ...
	\usepackage{../../hblib/hblibDatastructures} % lib data structures: \hString, \hStack, \hQueue
% ~~~~~~~~~~~~~~~~~~~~~~~~~~~~~~~~~~~~~~~~~~~~~~~~~~~~~~~~~~~~~~~~~ A hblib


\usepackage[T1]{fontenc}
\usepackage{lmodern}
\usepackage{booktabs}

\newcommand{\win}{\textbf{Windows}}
\newcommand{\mac}{\textbf{macOS}}

% long Java lines: wrap instead of running off the frame
\lstset{columns=fixed, breaklines=true, breakatwhitespace=false}

\SetAlCapFnt{\small}
\SetAlCapNameFnt{\small}
\SetKwProg{Fn}{function}{}{end}
\SetKw{KwError}{error}

% ---------------------------------------------------------------------
\title{Linked Lists}
\subtitle{Sophomore Level: Generic Singly and Doubly Linked Lists}
\author[Alpaydin]{Ruhan Alpaydin}

\institute[FCIS]{
    Faculty of Computer and Information Sciences\\
    Yeditepe University
}
\date{
	{\tiny \hbTimeStamp}
}

\begin{document}




% ~~~~~~~~~~~~~~~~~~~~~~~~~~~~~~~~~~~~~~ V frm
\begin{frame}[t,fragile]
  \titlepage
\end{frame}
% ~~~~~~~~~~~~~~~~~~~~~~~~~~~~~~~~~~~~~~ A frm




% ~~~~~~~~~~~~~~~~~~~~~~~~~~~~~~~~~~~~~~ V frm
\begin{frame}[t,fragile] \frametitle{Outline}
  \begin{columns}[T]
    \begin{column}{0.48\textwidth}
      \tableofcontents[sections={1-3}, hideallsubsections]
    \end{column}
    \begin{column}{0.48\textwidth}
      \tableofcontents[sections={4-6}, hideallsubsections]
    \end{column}
  \end{columns}
\end{frame}
% ~~~~~~~~~~~~~~~~~~~~~~~~~~~~~~~~~~~~~~ A frm

% ~~~~~~~~~~~~~~~~~~~~~~~~~~~~~~~~~~~~~~ V
\hSection[ADT]{The Interface}
% ~~~~~~~~~~~~~~~~~~~~~~~~~~~~~~~~~~~~~~ A




% ~~~~~~~~~~~~~~~~~~~~~~~~~~~~~~~~~~~~~~ V frm
\begin{frame}[t,fragile] \frametitle{\texttt{LinkedListADT<T>}}
  \lstcode[basicstyle=\ttfamily\scriptsize]{codeoutput/LinkedListADT.java}{Java}
  \begin{itemize}\small
    \item Seven methods, nothing more.
    \item \hCode{<T>} is the element type: any reference type.
    \item Both a singly linked and a doubly linked class can implement
          this interface: they have the same contract, different costs.
  \end{itemize}
\end{frame}
% ~~~~~~~~~~~~~~~~~~~~~~~~~~~~~~~~~~~~~~ A frm




% ~~~~~~~~~~~~~~~~~~~~~~~~~~~~~~~~~~~~~~ V frm
\begin{frame}[t,fragile] \frametitle{Why an interface first}
  \begin{columns}[T]
    \begin{column}{0.48\textwidth}
      \begin{block}{One contract}
        \small
        Both \hCode{SinglyLinked<T>} and \hCode{DoublyLinked<T>} obey
        \hCode{LinkedListADT<T>}.
        Any code written against the interface works with either.
      \end{block}
    \end{column}
    \begin{column}{0.48\textwidth}
      \begin{block}{Different costs}
        \small
        The contract is identical; only the cost of
        \hCode{removeLast} changes ($O(n)$ vs $O(1)$).
        That difference is the whole reason a doubly linked list exists.
      \end{block}
    \end{column}
  \end{columns}
  \vspace{0.4em}
  \small
  \hRemark{Same idea as the JDK}: \hCode{java.util.Queue} is an interface,
  \hCode{ArrayDeque} and \hCode{LinkedList} are two implementations.
\end{frame}
% ~~~~~~~~~~~~~~~~~~~~~~~~~~~~~~~~~~~~~~ A frm

% ~~~~~~~~~~~~~~~~~~~~~~~~~~~~~~~~~~~~~~ V
\hSection[Singly]{Singly Linked}
% ~~~~~~~~~~~~~~~~~~~~~~~~~~~~~~~~~~~~~~ A




% ~~~~~~~~~~~~~~~~~~~~~~~~~~~~~~~~~~~~~~ V frm
\begin{frame}[t,fragile] \frametitle{\texttt{SinglyLinked<T>}: fields and add/remove}
  \begin{columns}[T]
    \begin{column}{0.52\textwidth}
      \lstcode[basicstyle=\ttfamily\tiny, firstline=3, lastline=28]{codeoutput/SinglyLinked.java}{Java}
    \end{column}
    \begin{column}{0.44\textwidth}
      \lstcode[basicstyle=\ttfamily\tiny, firstline=30, lastline=40, firstnumber=30]{codeoutput/SinglyLinked.java}{Java}
      \lstterminal[basicstyle=\ttfamily\scriptsize]{codeoutput/SinglyLinked.txt}
    \end{column}
  \end{columns}
\end{frame}
% ~~~~~~~~~~~~~~~~~~~~~~~~~~~~~~~~~~~~~~ A frm




% ~~~~~~~~~~~~~~~~~~~~~~~~~~~~~~~~~~~~~~ V frm
\begin{frame}[t,fragile] \frametitle{\texttt{SinglyLinked<T>}: on paper}
  \begin{center}
  \begin{tikzpicture}
    \node at (-1.4,0) {\hCode{head}};
    \node[draw, minimum width=1.1cm, minimum height=0.55cm] (a) at (0,0) {2};
    \node[draw, minimum width=1.1cm, minimum height=0.55cm] (b) at (1.9,0) {1};
    \node[draw, minimum width=1.1cm, minimum height=0.55cm] (c) at (3.8,0) {3};
    \node at (5.3,0) {\hCode{null}};
    \draw[-latex] (-0.75,0) -- (a.west);
    \draw[-latex] (a.east)  -- (b.west);
    \draw[-latex] (b.east)  -- (c.west);
    \draw[-latex] (c.east)  -- (4.9,0);
    \node[font=\scriptsize] at (2.7,-0.65) {\hCode{tail} $\rightarrow$ the last node};
    \draw[-latex, thick, colorUsed] (2.7,-0.5) -- (3.8,-0.1);
  \end{tikzpicture}
  \end{center}
  \vspace{0.3em}
  \begin{columns}[T]
    \begin{column}{0.48\textwidth}
      \begin{block}{After the run}
        \small
        \hCode{contains 2: true}\\
        \hCode{2 1 3}\\
        \hCode{removeFirst: 2}\\
        \hCode{removeLast: 3}
      \end{block}
    \end{column}
    \begin{column}{0.48\textwidth}
      \begin{block}{Watch the costs}
        \small
        \hCode{addFirst} and \hCode{removeFirst} touch only \hCode{head}:
        $O(1)$.\\
        \hCode{removeLast} walks from \hCode{head} to find the node
        before \hCode{tail}: $O(n)$.
      \end{block}
    \end{column}
  \end{columns}
\end{frame}
% ~~~~~~~~~~~~~~~~~~~~~~~~~~~~~~~~~~~~~~ A frm

% ~~~~~~~~~~~~~~~~~~~~~~~~~~~~~~~~~~~~~~ V
\hSection[Doubly]{Doubly Linked}
% ~~~~~~~~~~~~~~~~~~~~~~~~~~~~~~~~~~~~~~ A




% ~~~~~~~~~~~~~~~~~~~~~~~~~~~~~~~~~~~~~~ V frm
\begin{frame}[t,fragile] \frametitle{\texttt{DoublyLinked<T>}: fields and node}
  \lstcode[basicstyle=\ttfamily\tiny, firstline=3, lastline=14, firstnumber=3]{codeoutput/DoublyLinked.java}{Java}
  \begin{itemize}\small
    \item The inner \hCode{Node<T>} has \hCode{item}, \hCode{next}, \hCode{prev}.
    \item \hCode{head} and \hCode{tail} are the two ends of the list.
    \item \hCode{size} keeps the count: no walking to get it.
    \item \hRemark{One extra reference per element}, in exchange for
          $O(1)$ on \emph{both} ends.
  \end{itemize}
\end{frame}
% ~~~~~~~~~~~~~~~~~~~~~~~~~~~~~~~~~~~~~~ A frm




% ~~~~~~~~~~~~~~~~~~~~~~~~~~~~~~~~~~~~~~ V frm
\begin{frame}[t,fragile] \frametitle{\texttt{DoublyLinked<T>}: \texttt{addFirst} and \texttt{addLast}}
  \lstcode[basicstyle=\ttfamily\tiny, firstline=16, lastline=32, firstnumber=16]{codeoutput/DoublyLinked.java}{Java}
  \begin{itemize}\small
    \item Both operations are $O(1)$: they touch only the two ends.
    \item Empty case: both \hCode{head} and \hCode{tail} point at the new node.
    \item The \hCode{prev} link is what makes \hCode{removeLast} possible later.
  \end{itemize}
\end{frame}
% ~~~~~~~~~~~~~~~~~~~~~~~~~~~~~~~~~~~~~~ A frm




% ~~~~~~~~~~~~~~~~~~~~~~~~~~~~~~~~~~~~~~ V frm
\begin{frame}[t,fragile] \frametitle{\texttt{DoublyLinked<T>}: \texttt{remove}}
  \lstcode[basicstyle=\ttfamily\tiny, firstline=34, lastline=52, firstnumber=34]{codeoutput/DoublyLinked.java}{Java}
  \begin{columns}[T]
    \begin{column}{0.48\textwidth}
      \begin{block}{The loop}
        \small
        Walk from \hCode{head} via \hCode{next}, compare \hCode{p.item}
        with the target.
      \end{block}
    \end{column}
    \begin{column}{0.48\textwidth}
      \begin{block}{The unlink}
        \small
        Two references per side:
        \hCode{p.prev.next = p.next} and
        \hCode{p.next.prev = p.prev}.
        The null checks handle the head and the tail.
      \end{block}
    \end{column}
  \end{columns}
\end{frame}
% ~~~~~~~~~~~~~~~~~~~~~~~~~~~~~~~~~~~~~~ A frm




% ~~~~~~~~~~~~~~~~~~~~~~~~~~~~~~~~~~~~~~ V frm
\begin{frame}[t,fragile] \frametitle{First run: \texttt{TestDoubly}}
  \begin{columns}[T]
    \begin{column}{0.54\textwidth}
      \lstcode[basicstyle=\ttfamily\tiny]{codeoutput/TestDoubly.java}{Java}
    \end{column}
    \begin{column}{0.42\textwidth}
      \lstterminal[basicstyle=\ttfamily\scriptsize]{codeoutput/TestDoubly.txt}
      \small
      \hCode{remove(2)} reports \hCode{true}; \hCode{remove(9)} would
      report \hCode{false}. The size drops by one when the remove succeeds.
    \end{column}
  \end{columns}
\end{frame}
% ~~~~~~~~~~~~~~~~~~~~~~~~~~~~~~~~~~~~~~ A frm

% ~~~~~~~~~~~~~~~~~~~~~~~~~~~~~~~~~~~~~~ V
\hSection[Iterator]{Iterator}
% ~~~~~~~~~~~~~~~~~~~~~~~~~~~~~~~~~~~~~~ A




% ~~~~~~~~~~~~~~~~~~~~~~~~~~~~~~~~~~~~~~ V frm
\begin{frame}[t,fragile] \frametitle{Making \texttt{SinglyLinked<T>} iterable}
  \lstcode[basicstyle=\ttfamily\tiny, firstline=42, lastline=52, firstnumber=42]{codeoutput/SinglyLinked.java}{Java}
  \begin{itemize}\small
    \item \hCode{SinglyLinked<T>} declares \hCode{implements Iterable<T>}.
    \item The anonymous inner class holds its own cursor \hCode{p},
          independent of the list's \hCode{head}.
    \item \hCode{hasNext()} is a null check; \hCode{next()} returns the
          item and advances \hCode{p}.
    \item This is the only way Java lets you write
          \hCode{for (T x : list)}.
  \end{itemize}
\end{frame}
% ~~~~~~~~~~~~~~~~~~~~~~~~~~~~~~~~~~~~~~ A frm




% ~~~~~~~~~~~~~~~~~~~~~~~~~~~~~~~~~~~~~~ V frm
\begin{frame}[t,fragile] \frametitle{Using the iterator}
  \begin{columns}[T]
    \begin{column}{0.54\textwidth}
      \lstcode[basicstyle=\ttfamily\tiny]{codeoutput/TestIterator.java}{Java}
    \end{column}
    \begin{column}{0.42\textwidth}
      \lstterminal[basicstyle=\ttfamily\scriptsize]{codeoutput/TestIterator.txt}
      \small
      The for-each loop calls \hCode{iterator()} once, then
      \hCode{hasNext()} / \hCode{next()} until it returns \hCode{false}.
      The list itself is not modified.
    \end{column}
  \end{columns}
\end{frame}
% ~~~~~~~~~~~~~~~~~~~~~~~~~~~~~~~~~~~~~~ A frm

% ~~~~~~~~~~~~~~~~~~~~~~~~~~~~~~~~~~~~~~ V
\hSection[Compare]{JDK Comparison}
% ~~~~~~~~~~~~~~~~~~~~~~~~~~~~~~~~~~~~~~ A




% ~~~~~~~~~~~~~~~~~~~~~~~~~~~~~~~~~~~~~~ V frm
\begin{frame}[t,fragile] \frametitle{Our class vs \texttt{java.util.LinkedList}}
  \begin{center}\footnotesize
  \begin{tabular}{L{0.22\textwidth}L{0.22\textwidth}L{0.46\textwidth}}
    \toprule
    \textbf{Your class} & \textbf{JDK} & \textbf{What the JDK adds} \\
    \midrule
    \hCode{SinglyLinked<T>}  & (none directly) & --- \\
    \hCode{DoublyLinked<T>}  & \hCode{LinkedList<T>}
        & \hCode{List}, \hCode{Deque}, \hCode{Queue}, \hCode{Iterator},
          \hCode{ListIterator}, \hCode{Cloneable}, \hCode{Serializable} \\
    \bottomrule
  \end{tabular}
  \end{center}
  \vspace{0.4em}
  \small
  \hRemark{The JDK does not ship a singly linked list class}:
  the doubly linked \hCode{LinkedList} covers every use that a singly
  linked list would, at the cost of one extra reference per element.
\end{frame}
% ~~~~~~~~~~~~~~~~~~~~~~~~~~~~~~~~~~~~~~ A frm




% ~~~~~~~~~~~~~~~~~~~~~~~~~~~~~~~~~~~~~~ V frm
\begin{frame}[t,fragile] \frametitle{The same list, two vocabularies}
  \begin{center}\footnotesize
  \begin{tabular}{@{}lll@{}}
    \toprule
    \textbf{What you want} & \textbf{In your class} & \textbf{In \texttt{java.util.LinkedList}} \\
    \midrule
    add at head          & \hCode{addFirst(x)}    & \hCode{addFirst(x)} \\
    add at tail          & \hCode{addLast(x)}     & \hCode{addLast(x)} \\
    remove at head       & \hCode{removeFirst()}  & \hCode{removeFirst()} \\
    remove at tail       & \hCode{removeLast()}   & \hCode{removeLast()} \\
    find an element      & \hCode{contains(x)}    & \hCode{contains(x)} \\
    number of elements   & \hCode{size()}         & \hCode{size()} \\
    walk all elements    & for-each               & for-each, \hCode{descendingIterator()} \\
    \bottomrule
  \end{tabular}
  \end{center}
  \vspace{0.3em}
  \small
  \hRemark{Same interface, same names, same $O(1)$ costs} for the ends.
\end{frame}
% ~~~~~~~~~~~~~~~~~~~~~~~~~~~~~~~~~~~~~~ A frm

% ~~~~~~~~~~~~~~~~~~~~~~~~~~~~~~~~~~~~~~ V
\hSection[Wrap]{Wrap-up}
% ~~~~~~~~~~~~~~~~~~~~~~~~~~~~~~~~~~~~~~ A




% ~~~~~~~~~~~~~~~~~~~~~~~~~~~~~~~~~~~~~~ V frm
\begin{frame}[t,fragile] \frametitle{Complexities}
  \begin{center}\small
  \begin{tabular}{@{}lccc@{}}
    \toprule
    \textbf{Operation} & \textbf{Singly} & \textbf{Doubly} & \textbf{Why} \\
    \midrule
    \hCode{addFirst}     & $O(1)$ & $O(1)$ & touches only \hCode{head} \\
    \hCode{removeFirst}  & $O(1)$ & $O(1)$ & touches only \hCode{head} \\
    \hCode{addLast}      & $O(1)$ & $O(1)$ & \hCode{tail} is kept \\
    \hCode{removeLast}   & $O(n)$ & $O(1)$ & singly must walk; doubly uses \hCode{prev} \\
    \hCode{contains}     & $O(n)$ & $O(n)$ & no jump table in either \\
    \hCode{remove(x)}    & $O(n)$ & $O(n)$ & must find \hCode{x} first \\
    memory per element   & 2 refs & 3 refs & \hCode{item} and \hCode{next} vs \hCode{item}, \hCode{next}, \hCode{prev} \\
    \bottomrule
  \end{tabular}
  \end{center}
  \vspace{0.3em}
  \small
  \hRemark{\texttt{removeLast} is the only row where the two classes differ}.
\end{frame}
% ~~~~~~~~~~~~~~~~~~~~~~~~~~~~~~~~~~~~~~ A frm




% ~~~~~~~~~~~~~~~~~~~~~~~~~~~~~~~~~~~~~~ V frm
\begin{frame}[t,fragile] \frametitle{Big ideas}
  \begin{itemize}
    \item A singly linked node holds \hCode{item} and \hCode{next};
          a doubly linked node also holds \hCode{prev}.
    \item Adding and removing at the \textbf{head} is $O(1)$ in both.
    \item \hCode{removeLast} on a singly linked list is $O(n)$;
          on a doubly linked list it is $O(1)$ because of \hCode{prev}.
    \item The \hCode{Node} class is \hCode{private static}; it exists
          only to serve the list.
    \item An iterator is a cursor over the chain; it is the only way to
          walk a linked list with the for-each syntax in Java.
    \item The JDK's \hCode{LinkedList} is a doubly linked list that
          also implements \hCode{List}, \hCode{Deque}, \hCode{Queue}.
  \end{itemize}
\end{frame}
% ~~~~~~~~~~~~~~~~~~~~~~~~~~~~~~~~~~~~~~ A frm




% ~~~~~~~~~~~~~~~~~~~~~~~~~~~~~~~~~~~~~~ V frm
\begin{frame}[t,fragile] \frametitle{Check yourself}
  \begin{enumerate}\small
    \item Why must the \hCode{Node} class be \hCode{static} when it is
          nested inside a generic class?
    \item What does \hCode{removeLast} cost on \hCode{SinglyLinked}
          and why?
    \item Why does \hCode{addLast} need a \hCode{tail} reference on both classes?
    \item What does the iterator hold between two calls to \hCode{next()}?
    \item What is the difference between \hCode{size()} and the number
          of nodes reachable from \hCode{head}? (When would they differ?)
    \item When would you choose \hCode{LinkedList} over \hCode{ArrayList}
          in the JDK?
  \end{enumerate}
\end{frame}
% ~~~~~~~~~~~~~~~~~~~~~~~~~~~~~~~~~~~~~~ A frm




% ~~~~~~~~~~~~~~~~~~~~~~~~~~~~~~~~~~~~~~ V frm
\begin{frame}[t,fragile] \frametitle{Tests (Not run here)}
  \begin{columns}[T]
    \begin{column}{0.56\textwidth}
      \begin{block}{JUnit: \emph{Not run here}}
        \small
        \begin{lstlisting}[language=Java, basicstyle=\ttfamily\tiny, frame=single]
@Test
void addFirstGrowsSize() {
    SinglyLinked<Integer> s = new SinglyLinked<>();
    s.addFirst(1); s.addFirst(2);
    assertEquals(2, s.size());
    assertTrue(s.contains(1));
    assertEquals(2, s.removeFirst());
    assertEquals(1, s.size());
}
        \end{lstlisting}
      \end{block}
    \end{column}
    \begin{column}{0.4\textwidth}
      \begin{block}{What it checks}
        \small
        Each \hCode{@Test} is a small program that \emph{asserts} the
        contract: size goes up by one, the removed value is the one
        you inserted first, and \hCode{contains} agrees.
      \end{block}
    \end{column}
  \end{columns}
  \vspace{0.3em}
  \small
  \hRemark{Run it once, and the assertions become a contract you can trust}.
\end{frame}
% ~~~~~~~~~~~~~~~~~~~~~~~~~~~~~~~~~~~~~~ A frm




% ~~~~~~~~~~~~~~~~~~~~~~~~~~~~~~~~~~~~~~ V frm
\begin{frame}[t,fragile] \frametitle{Practice}
  \begin{enumerate}\small
    \item Add \hCode{get(int i)} to \hCode{SinglyLinked<T>}. What is its
          complexity, and why is it not used in the tests?
    \item Add a \hCode{toString()} to both classes so
          \hCode{System.out.println(s)} prints the elements in order.
    \item Add \hCode{removeFirst} to \hCode{SinglyLinked<T>} if it is missing.
          Then use it as a stack in a small program.
    \item Add \hCode{remove(int index)} to \hCode{DoublyLinked<T>}.
          Walk from \hCode{head} or from \hCode{tail}, whichever is closer.
    \item Make \hCode{DoublyLinked<T>} implement \hCode{Iterable<T>},
          both forward and backward.
    \item Write a JUnit test that verifies the complexity table: add one
          million elements, call \hCode{removeLast}, and show that the
          doubly linked version does not slow down.
  \end{enumerate}
\end{frame}
% ~~~~~~~~~~~~~~~~~~~~~~~~~~~~~~~~~~~~~~ A frm




% ~~~~~~~~~~~~~~~~~~~~~~~~~~~~~~~~~~~~~~ V frm
\begin{frame}[t,fragile] \frametitle{References}
  \begin{enumerate}\small
    \item Robert Lafore, \emph{Data Structures \& Algorithms in Java},
          2nd ed., Sams, 2002. (The \hCode{Link}/\hCode{LinkList}
          style reused in the freshman deck.)
    \item Oracle, \emph{Java SE 21 API Specification}:
          \hCode{java.util.LinkedList},
          \hCode{java.util.List},
          \hCode{java.util.Deque},
          \hCode{java.util.Iterator}.
    \item Oracle, \emph{The Java Tutorials: Trail: Collections}.
    \item Joshua Bloch, \emph{Effective Java}, 3rd ed., Addison-Wesley,
          2018. (Items on \hCode{equals}, \hCode{hashCode}, \hCode{toString}.)
    \item Herbert Schildt, \emph{Java: The Complete Reference},
          12th ed., McGraw-Hill, 2021.
  \end{enumerate}
\end{frame}
% ~~~~~~~~~~~~~~~~~~~~~~~~~~~~~~~~~~~~~~ A frm




% ~~~~~~~~~~~~~~~~~~~~~~~~~~~~~~~~~~~~~~ V frm
\begin{frame}[t,fragile]
  \begin{center}
    \vspace{4cm}
    \hRemark{\Huge Questions?}
  \end{center}
\end{frame}
% ~~~~~~~~~~~~~~~~~~~~~~~~~~~~~~~~~~~~~~ A frm




\end{document}